\chapter{Banking Sector \& Institution Profiles}
\label{ch:banking}

This chapter provides a comprehensive, viva-oriented guide to Bangladesh's banking sector. It includes tactical justification for bank choices (specifically defending Sonali Bank PLC as the premier choice), detailed institutional profiles of state-owned commercial and specialized financial institutions, fundamental banking terms, regulatory frameworks, and deep technical coverage of banking IT and payment architectures.

\medskip

\section{First-Choice Defense: Sonali Bank PLC}

\subsection{Why is Sonali Bank PLC Your First Choice?}

\begin{quote}
\textit{Model Viva Answer for an IT / Engineering Candidate:}
\end{quote}

``Sonali Bank PLC is my first choice for three core reasons: \textbf{national impact}, \textbf{unmatched IT infrastructure scale}, and \textbf{institutional stability}.

First, Sonali Bank is the largest state-owned commercial bank in Bangladesh, operating a massive nationwide network of \textbf{1,235 branches} and serving as the financial lifeline of the government. Beyond standard commercial operations, it acts as an agent of Bangladesh Bank, executing national treasury functions, revenue collection through the Sonali Payment Gateway, social safety-net disbursements, and student stipends. Working here allows an IT professional to directly contribute to the backbone of Bangladesh's digital economy.

Second, from a Computer Science and IT perspective, Sonali Bank offers an engineering challenge that no other financial institution in the country can match. Managing the Core Banking Solution (CBS) for over a thousand branches, handling terabytes of mission-critical financial data, supporting high-throughput transactions through \textbf{Sonali eSheba}, \textbf{Sonali eWallet}, and \textbf{BLAZE} instant remittance, and securing integrations with Bangladesh Bank's RTGS, NPSB, and BEFTN represents a tier-one distributed computing ecosystem.

Third, Sonali Bank is a 100\% government-owned public limited company with sovereign backing, excellent career progression up to senior executive management, and structured training through the Sonali Bank Staff College. It offers the ideal balance of public service, career prestige, and deep technological growth.''

\medskip

\subsection{Sonali Bank PLC: Institutional Factsheet}

\begin{itemize}
  \item \textbf{Genesis \& Legal Foundation:} Established on March 26, 1972, under the Bangladesh Banks (Nationalisation) Order, 1972 (Presidential Order No. 26 of 1972), by amalgamating the former National Bank of Pakistan, Premier Bank, and Bank of Bahawalpur.
  \item \textbf{Corporatization:} Reconstituted as a Public Limited Company on \textbf{November 15, 2007}, with 100\% government shareholding, operating under the Companies Act, 1994. Rebranded to \textbf{Sonali Bank PLC} (formerly Sonali Bank Limited) in accordance with the amended company legislation.
  \item \textbf{Head Office:} Motijheel Commercial Area, Dhaka-1000, Bangladesh.
  \item \textbf{Slogan:} \textit{Bishshosto o Smart} (Trustworthy \& Smart).
  \item \textbf{Capital Structure:}
  \begin{itemize}
    \item \textbf{Authorized Capital:} Tk. 6,000.00 Crore.
    \item \textbf{Paid-up Capital:} Tk. 4,530.00 Crore.
  \end{itemize}
  \item \textbf{Branch Network (Total: 1,235):}
  \begin{itemize}
    \item \textbf{Domestic Branches:} 1,233 (Rural: 697; Urban: 536).
    \item \textbf{Overseas Branches:} 2 full-fledged branches in India (Kolkata and Siliguri).
    \item \textbf{Administrative Tiers:} 18 General Manager (GM) Offices and 69 Principal Offices supervising operations nationwide.
  \end{itemize}
  \item \textbf{Subsidiaries \& Global Presence:}
  \begin{itemize}
    \item \textbf{Sonali Exchange Company Inc. (SECI):} Operating 10 branches across the United States for foreign remittance.
    \item \textbf{Sonali Bank (UK) Limited:} Operating 2 branches in the United Kingdom, alongside \textbf{Sonali Pay UK Limited}.
    \item \textbf{Sonali Investment Limited (SIL):} Merchant banking subsidiary operating 5 branches (4 in Dhaka, 1 in Khulna).
    \item \textbf{Representative Offices:} 3 foreign representative offices in Jeddah (KSA), Riyadh (KSA), and Kuwait.
    \item \textbf{Joint Venture Associate:} \textbf{Sonali Intellect Limited} (a prominent IT joint venture focused on financial software solutions).
  \end{itemize}
  \item \textbf{Amalgamation / Merger Milestone:} Under a Memorandum of Understanding (MoU) signed on \textbf{May 12, 2024}, Bangladesh Development Bank PLC (BDBL) is undergoing a voluntary merger process into Sonali Bank PLC as part of the national banking sector rationalization roadmap.
\end{itemize}

\medskip

\subsection{Managerial Structure and Corporate Governance}

\begin{itemize}
  \item \textbf{Board of Directors:} The supreme policy-making body, comprising \textbf{11 directors} appointed by the Government of Bangladesh, consisting of senior civil servants, financial experts, and independent directors.
  \item \textbf{Current Chairman:} \textbf{Muhammed Mahbubur Rahman} (Chairman, Current Charge).
  \item \textbf{Current Managing Director \& CEO:} \textbf{Md. Shawkat Ali Khan} (heads executive administration and operational management).
  \item \textbf{Executive Hierarchy (Top-Down):}
\end{itemize}

\begin{lstlisting}
Board of Directors
       |
Managing Director & CEO
       |
Deputy Managing Directors (DMD)
       |
General Managers (GM)  [Circle / Divisional Heads]
       |
Deputy General Managers (DGM) [Department / Principal Office Heads]
       |
Assistant General Managers (AGM)
       |
Senior Principal Officers (SPO)
       |
Principal Officers (PO)
       |
Senior Officers (SO) / Assistant Programmers (AP) / SO (IT)
       |
Officers / Junior Officers
\end{lstlisting}

\medskip

\subsection{Special Government \& Agency Functions of Sonali Bank}

Sonali Bank is unique among commercial banks because it acts as an \textbf{agent of Bangladesh Bank}:

\begin{enumerate}
  \item \textbf{Treasury \& Sub-Treasury Operations:} Where Bangladesh Bank lacks physical branch presence, Sonali Bank chest branches handle government revenue receipts, currency storage, and payments on behalf of the central bank.
  \item \textbf{Utility \& Government Fee Collection:} Automated collection of NBR customs duty, income tax, VAT via NBR-Sonali e-Payment, e-challan, passport fees, road transport (BRTA) fees, land registration fees, and university entrance fees.
  \item \textbf{Social Safety Net \& Public Disbursements:} Direct distribution of government freedom fighter allowances, elderly pensions, widow allowances, rural development micro-credits, and agricultural subsidies.
  \item \textbf{National Savings Certificates (Sanchayapatra):} The largest issuer and settlement agent for national savings schemes and prize bonds.
  \item \textbf{Foreign Food \& Commodity Aid:} Facilitating government import letters of credit (LC) for petroleum (BPC), food grains, fertilizers, and power equipment.
\end{enumerate}

\medskip

\subsection{Sonali Bank's Digital Ecosystem \& IT Initiatives}

For an IT candidate, being familiar with Sonali Bank's software infrastructure is essential:

\begin{itemize}
  \item \textbf{Core Banking Solution (CBS):} Centralized database linking all 1,235 domestic and foreign branches in real time.
  \item \textbf{Sonali eSheba:} Mobile-first digital onboarding platform enabling account opening through biometric e-KYC in under five minutes.
  \item \textbf{Sonali eWallet:} Smartphone digital wallet supporting QR-code payments, mobile recharges, fund transfers, and bill payments.
  \item \textbf{Sonali Payment Gateway (SPG):} Dedicated online payment aggregator powering automated revenue collection for over 50 government ministries, autonomous universities, and statutory boards.
  \item \textbf{BLAZE Service:} 24/7 instant cross-border inward remittance settlement service, crediting beneficiary bank accounts across Bangladesh within five seconds.
  \item \textbf{Binimoy Integration:} Seamless interoperability with the Interoperable Digital Transaction Platform (IDTP), bridging accounts between banks, Mobile Financial Services (bKash, Nagad, Rocket), and payment service providers.
\end{itemize}

\medskip

\section{Key Financial Institutions of Bangladesh}

\subsection{Bangladesh Bank (The Central Bank)}

\begin{itemize}
  \item \textbf{Genesis:} Established on \textbf{December 16, 1971}, pursuant to the Bangladesh Bank Order, 1972 (Presidential Order No. 127 of 1972), taking over the Dhaka branch of the State Bank of Pakistan.
  \item \textbf{Governor:} \textbf{Md. Mostaqur Rahman} assumed office on \textbf{February 26, 2026}, as the \textbf{14th Governor} of Bangladesh Bank (succeeding Dr. Ahsan H. Mansur).
  \item \textbf{First Governor:} A. N. M. Hamidullah (1972--1974).
  \item \textbf{Branch Offices (10):} Motijheel (Head Office), Sadarghat, Chattogram, Khulna, Rajshahi, Bogura, Sylhet, Barishal, Rangpur, and Mymensingh.
  \item \textbf{Core Mandates:}
  \begin{enumerate}
    \item Formulating and implementing monetary and credit policies to maintain price stability.
    \item Regulating and supervising scheduled banks and Non-Bank Financial Institutions (NBFIs).
    \item Holding and managing the official foreign exchange reserves of Bangladesh.
    \item Sole authority to issue currency notes: Central bank notes include \textbf{Tk. 10, 20, 50, 100, 200, 500, and 1000} (signed by the Governor).
    \item Government notes and coins include \textbf{Tk. 1, 2, and 5} (issued by the Ministry of Finance, signed by the Finance Secretary).
    \item Operator of national payment switches (BACH, RTGS, NPSB, BEFTN).
    \item Acting as banker to the government and lender of last resort to commercial banks.
  \end{enumerate}
\end{itemize}

\medskip

\subsection{State-Owned Commercial Banks (SOCBs)}

Bangladesh has four major state-owned commercial banks:

\begin{enumerate}
  \item \textbf{Sonali Bank PLC:} Largest bank, 1,235 branches, comprehensive treasury agent.
  \item \textbf{Agrani Bank PLC:}
  \begin{itemize}
    \item Established under PO 26 of 1972 by merging Habib Bank and Commerce Bank; corporatized in 2007.
    \item Chairman: Syed Abu Naser Bukhtear Ahmed; Managing Director (Current Charge): Parsoma Alam.
    \item Branch Network: \textbf{979 branches} across 11 circle offices; pioneer in overseas remittance operations through Agrani Exchange houses.
  \end{itemize}
  \item \textbf{Janata Bank PLC:}
  \begin{itemize}
    \item Established under PO 26 of 1972 by merging United Bank and Union Bank; corporatized in 2007.
    \item Second-largest state-owned commercial bank with approximately \textbf{920+ branches} and 4 foreign branches in the UAE (Abu Dhabi, Dubai, Sharjah, Al-Ain).
    \item Plays a major role in import-export trade finance and industrial term financing.
  \end{itemize}
  \item \textbf{Rupali Bank PLC:}
  \begin{itemize}
    \item Established under PO 26 of 1972 by amalgamating Muslim Commercial Bank, Australasia Bank, and Standard Bank.
    \item Unique distinction: The \textbf{only publicly listed state-owned commercial bank} on the Dhaka Stock Exchange (DSE) and Chittagong Stock Exchange (CSE), with government holding approximately 90\%+ shares.
    \item Network: Around 585 branches.
  \end{itemize}
\end{enumerate}

\medskip

\subsection{Specialized Banks \& Development Financial Institutions}

Specialized banks are created under specific statutory acts to fulfill targeted socio-economic developmental objectives:

\begin{enumerate}
  \item \textbf{Bangladesh Krishi Bank (BKB):}
  \begin{itemize}
    \item Established under the Bangladesh Krishi Bank Order, 1973 (Presidential Order No. 27 of 1973).
    \item Chairman: Mohammad Nurul Amin; Managing Director: Sanchia Binte Ali.
    \item Network: Over \textbf{1,030 branches}, the vast majority situated in rural areas.
    \item Purpose: Providing agricultural credit, financing crop production, livestock, fisheries, agro-machinery, and rural cottage industries.
  \end{itemize}
  \item \textbf{Rajshahi Krishi Unnayan Bank (RAKUB):}
  \begin{itemize}
    \item Established on March 15, 1986, pursuant to President's Ordinance No. 58 of 1986, taking over BKB's branches in Rajshahi and Rangpur administrative divisions.
    \item Headquarters: Rajshahi; operates around 383 branches across 16 northern districts to boost the agricultural economy of the northwest grain belt.
  \end{itemize}
  \item \textbf{Bangladesh Development Bank PLC (BDBL):}
  \begin{itemize}
    \item Formed on November 16, 2009, by merging two historic development institutions: \textbf{Bangladesh Shilpa Bank (BSB)} and \textbf{Bangladesh Shilpa Rin Sangstha (BSRS)}.
    \item Network: 51 branches specializing in industrial financing and corporate loans.
    \item Status: Currently undergoing an approved merger into Sonali Bank PLC.
  \end{itemize}
  \item \textbf{Investment Corporation of Bangladesh (ICB):}
  \begin{itemize}
    \item Established on October 1, 1976, under the Investment Corporation of Bangladesh Ordinance, 1976 (now governed by the ICB Act, 2014).
    \item Chairman: Professor Abu Ahmed; Managing Director: Niranjan Chandra Debnath.
    \item Role: The apex national non-bank capital market institution. Functions include underwriting shares, managing mutual funds, acting as debenture trustee, stock market stabilization, and investor portfolio management.
    \item Subsidiaries: ICB Capital Management Ltd., ICB Asset Management Company Ltd., and ICB Securities Trading Company Ltd.
  \end{itemize}
  \item \textbf{Probashi Kollyan Bank (PKB):}
  \begin{itemize}
    \item Established under the Probashi Kollyan Bank Act, 2010.
    \item Purpose: Specialized bank providing collateral-free rehabilitation loans, migration loans for overseas employment candidates, and financial security to migrant wage earners.
  \end{itemize}
  \item \textbf{Karmasangsthan Bank:}
  \begin{itemize}
    \item Established under the Karmasangsthan Bank Act, 1998.
    \item Purpose: Financing unemployed youths with credit, entrepreneurship training, and seed capital for micro-enterprises to generate self-employment.
  \end{itemize}
\end{enumerate}

\medskip

\section{Essential Banking Terms \& Viva Concepts}

\subsection{Central Banking \& Monetary Policy Tools}

\begin{itemize}
  \item \textbf{Cash Reserve Ratio (CRR):} The percentage of a commercial bank's total Demand and Time Liabilities (DTL) that must be maintained in \textbf{cash with Bangladesh Bank}. CRR pays no interest and controls excess liquidity directly.
  \item \textbf{Statutory Liquidity Ratio (SLR):} The percentage of net demand and time liabilities that a bank must maintain in \textbf{liquid assets} (cash, gold, unencumbered approved government securities, treasury bills/bonds). For conventional commercial banks in Bangladesh, the statutory requirement is 13\% SLR + 4.0\% CRR.
  \item \textbf{Repo (Repurchase Agreement):} Bangladesh Bank lends short-term funds to commercial banks against government securities to inject liquidity into the banking system.
  \item \textbf{Reverse Repo / Standing Deposit Facility (SDF):} The rate at which commercial banks park their surplus liquidity with Bangladesh Bank, withdrawing money from circulation.
  \item \textbf{Policy Rate:} The primary benchmark interest rate (formerly Repo rate) set by Bangladesh Bank's Monetary Policy Statement to signal the direction of borrowing costs and fight inflation.
  \item \textbf{Bank Rate:} The long-term statutory interest rate at which Bangladesh Bank extends credit to scheduled banks and the government.
  \item \textbf{Inflation (Headline vs. Core):} \textit{Headline Inflation} measures total consumer price index (CPI) changes including volatile food and energy components; \textit{Core Inflation} strips out food and energy volatility to capture structural price pressures.
\end{itemize}

\medskip

\subsection{Bank Account Types \& Operational Instruments}

\begin{itemize}
  \item \textbf{Current Account (CA):} Running account designed for businesses and trade; unlimited deposits and withdrawals during banking hours; traditionally pays zero interest; allows overdraft facilities.
  \item \textbf{Savings Bank Account (SB):} Deposit account meant for individuals and non-profit institutions to encourage savings; earns interest; subject to daily or weekly withdrawal restrictions without notice.
  \item \textbf{Special Notice Deposit (SND):} Deposit account that earns interest, but funds can only be withdrawn upon giving prior notice (usually 7 or 30 days) to the bank.
  \item \textbf{Fixed Deposit Receipt (FDR):} Term deposit account where funds are locked for a specified tenure (e.g., 3 months, 1 year, 3 years) at a fixed pre-agreed interest rate; premature encashment incurs interest penalties.
  \item \textbf{Recurring Deposit (RD) / DPS (Deposit Pension Scheme):} Monthly installment-based savings scheme maturing into a lump-sum amount with compound interest.
  \item \textbf{Foreign Currency Accounts:}
  \begin{itemize}
    \item \textbf{NFCA (Non-Resident Foreign Currency Account):} Held by Bangladeshi nationals residing and earning abroad; tax-free interest; freely repatriable abroad.
    \item \textbf{RFCD (Resident Foreign Currency Deposit):} Allowed for resident Bangladeshis returning from abroad with unspent foreign exchange.
  \end{itemize}
\end{itemize}

\medskip

\subsection{Cheques \& The Negotiable Instruments Act, 1881}

\begin{itemize}
  \item \textbf{Negotiable Instrument:} A signed, transferable document that promises unconditional payment of a specific sum of money on demand or at a set future date (e.g., Cheques, Promissory Notes, Bills of Exchange).
  \item \textbf{Bearer Cheque:} Payable to whoever possesses or presents the cheque over the bank counter; carries maximum fraud risk.
  \item \textbf{Order Cheque:} Payable only to the person named on the cheque or to their order upon endorsement and verification of identity.
  \item \textbf{Crossed Cheque:} Two parallel transverse lines drawn across the upper left corner. Cannot be encashed over the counter; must be collected through a bank account:
  \begin{itemize}
    \item \textbf{General Crossing:} Two lines containing ``\& Co.'' or no text.
    \item \textbf{Special Crossing:} Includes the name of a specific collecting bank.
    \item \textbf{Account Payee Crossing (A/C Payee Only):} Strictly limits credit to the payee's account; destroys negotiability and prevents further endorsement.
  \end{itemize}
  \item \textbf{Stale Cheque:} A cheque presented for payment more than \textbf{6 months} after its date of issuance (cannot be honored without re-validation by the drawer).
  \item \textbf{Post-Dated Cheque (PDC):} A cheque bearing a date later than the date on which it was handed over; cannot be paid before the date arrives.
  \item \textbf{Section 138 of NI Act, 1881 (Dishonour of Cheque):} When a cheque bounces due to insufficient funds, the drawer commits a criminal offense punishable by imprisonment up to one year, or a fine up to three times the cheque amount, or both, subject to statutory 30-day legal notice.
\end{itemize}

\medskip

\subsection{Lending, Securities \& Credit Terminology}

\begin{itemize}
  \item \textbf{Working Capital Loan:} Short-term credit facility to finance day-to-day business operations (inventory, wages, accounts receivable).
  \item \textbf{Term Loan:} Medium or long-term loan for purchasing capital assets, machinery, plant expansion, or infrastructure, repayable in periodic installments.
  \item \textbf{Pledge:} Delivery of physical possession of goods/chattel to the bank as security for a loan (e.g., gold loan, warehouse commodities). The bank holds possession while the borrower retains ownership.
  \item \textbf{Hypothecation:} The borrower retains both possession and ownership of movable goods (raw materials, store inventory, vehicles), but creates an equitable charge/hypothecation deed in favor of the bank.
  \item \textbf{Mortgage:} Legal transfer of an interest in specific immovable property (land, building, factory) for securing payment of money advanced (governed by the Transfer of Property Act, 1882).
  \item \textbf{Lien:} The legal right of the bank to retain custody of property or securities belonging to the borrower until the underlying debt is satisfied (Banker's General Lien under Section 171 of Contract Act, 1872).
\end{itemize}

\medskip

\subsection{Non-Performing Loans (NPL) \& Recovery}

\begin{itemize}
  \item \textbf{Non-Performing Loan (NPL):} A loan where scheduled interest or principal repayments remain overdue beyond regulatory forbearance windows.
  \item \textbf{Classification Categories (Bangladesh Bank Master Circular):}
  \begin{enumerate}
    \item \textbf{Special Mention Account (SMA):} Overdue for 30 to 89 days (early warning signal).
    \item \textbf{Sub-Standard (SS):} Overdue for 90 to 179 days (3 months to under 6 months). Requires 20\% provisioning.
    \item \textbf{Doubtful (DF):} Overdue for 180 to 269 days (6 months to under 9/12 months). Requires 50\% provisioning.
    \item \textbf{Bad/Loss (B/L):} Overdue for 270/360 days or more. High probability of default. Requires 100\% provisioning.
  \end{enumerate}
  \item \textbf{Rescheduling:} Extending the loan repayment period or altering installment amounts when a borrower faces genuine distress, subject to mandatory down payment.
  \item \textbf{Restructuring:} Changing multiple terms of the credit contract (interest rates, grace periods, principal capitalization) for large viable corporate debts.
  \item \textbf{Write-off:} Removing a bad/loss loan from the active balance sheet after maintaining 100\% provisions and filing a money suit in the Artha Rin Adalat, though legal recovery efforts continue.
  \item \textbf{Artha Rin Adalat Ain, 2003:} Specialized Money Loan Court Act designed for accelerated judicial recovery of bad debts and foreclosure of mortgaged properties.
\end{itemize}

\medskip

\subsection{Prudential Standards: Basel III \& CAMELS}

\begin{itemize}
  \item \textbf{Basel III Accords:} International regulatory framework developed by the Basel Committee on Banking Supervision (BCBS) to strengthen capital adequacy, stress testing, and market liquidity risk.
  \item \textbf{Capital Adequacy Ratio (CAR / CRAR):} Measures a bank's available capital against its total Risk-Weighted Assets (RWA):
  \[
  \text{CRAR} = \frac{\text{Tier 1 Capital} + \text{Tier 2 Capital}}{\text{Total Risk-Weighted Assets (Credit + Market + Operational Risk)}} \times 100\%
  \]
  Bangladesh Bank mandates a minimum CRAR of \textbf{10\%}, plus an additional \textbf{2.5\% Capital Conservation Buffer (CCB)}, totaling 12.5\%.
  \item \textbf{Tier 1 (Core Capital):} Highest quality permanent capital, including Common Equity Tier 1 (paid-up capital, statutory reserve, retained earnings).
  \item \textbf{Tier 2 (Supplementary Capital):} Subordinated debt, general provisions against unclassified loans, and revaluation reserves.
  \item \textbf{CAMELS Rating:} Supervisory rating system used by Bangladesh Bank during annual inspections:
  \begin{itemize}
    \item \textbf{C} --- Capital Adequacy
    \item \textbf{A} --- Asset Quality (NPL ratio, loan concentration)
    \item \textbf{M} --- Management Soundness
    \item \textbf{E} --- Earnings \& Profitability (ROA, ROE)
    \item \textbf{L} --- Liquidity (ADR/LDR, LCR, NSFR)
    \item \textbf{S} --- Sensitivity to Market Risk (interest rate and FX risk)
  \end{itemize}
\end{itemize}

\medskip

\subsection{Foreign Exchange, Remittance \& Trade Finance}

\begin{itemize}
  \item \textbf{Letter of Credit (LC):} A binding commitment issued by an importer's bank (Issuing Bank) promising to pay an agreed sum to the exporter (Beneficiary) upon presentation of strictly complying shipping and trade documents.
  \item \textbf{Key Parties to an LC:}
  \begin{enumerate}
    \item \textbf{Applicant / Importer:} Requests the bank to open the credit.
    \item \textbf{Issuing Bank:} Opens the LC and assumes primary payment liability.
    \item \textbf{Advising Bank:} Bank in the exporter's country that authenticates and delivers the LC to the exporter.
    \item \textbf{Beneficiary / Exporter:} The seller in whose favor the credit is established.
    \item \textbf{Confirming Bank:} Adds a definitive guarantee to pay if the issuing bank defaults.
    \item \textbf{Negotiating Bank:} Purchases or discounts the drafts from the beneficiary.
  \end{enumerate}
  \item \textbf{Back-to-Back LC:} An LC opened by a garment manufacturer or exporter using their export LC as security to import raw materials (fabrics, accessories) without upfront cash.
  \item \textbf{Sight LC vs. Usance LC:} A \textit{Sight LC} requires payment immediately upon presentation of documents; a \textit{Usance (Deferred) LC} allows payment after a specified credit period (e.g., 60, 90, or 180 days).
  \item \textbf{Nostro vs. Vostro Accounts:}
  \begin{itemize}
    \item \textbf{Nostro Account} (``Our money at your bank''): An account maintained by a domestic bank (e.g., Sonali Bank) in foreign currency with an overseas bank (e.g., Citibank NY in US Dollars).
    \item \textbf{Vostro Account} (``Your money at our bank''): An account held by a foreign bank with a domestic bank in local currency (Taka).
    \item \textbf{Loro Account}: A third-party bank referencing another bank's Nostro/Vostro account.
  \end{itemize}
\end{itemize}

\medskip

\section{Information Technology \& Fintech in Banking}

\subsection{Core Banking Solutions (CBS)}

The Core Banking Solution is the centralized back-end database and transaction processing engine that handles real-time accounting across all connected branches:

\begin{itemize}
  \item \textbf{Architecture:} Multi-tier architecture consisting of Presentation layer (branch terminals, internet banking, mobile apps), Application/Business Logic layer (accounting rules, loan interest calculators, limits), and Database tier (clustered relational engines like Oracle, IBM DB2, or PostgreSQL).
  \item \textbf{ACID Guarantees:} Ensures strict \textbf{Atomicity, Consistency, Isolation, and Durability} across distributed concurrent financial transactions.
  \item \textbf{High Availability \& Disaster Recovery (DR):}
  \begin{itemize}
    \item \textbf{DC (Data Center):} Primary site processing active transactions.
    \item \textbf{DR (Disaster Recovery Site):} Geographically separated replica configured in \textbf{Active-Passive} (synchronous or near-synchronous log shipping) or \textbf{Active-Active} load-balanced topology.
    \item \textbf{RPO (Recovery Point Objective):} Maximum acceptable data loss duration during a disaster (ideally near zero for financial ledgers).
    \item \textbf{RTO (Recovery Time Objective):} Maximum acceptable downtime to switch operations to the DR site.
  \end{itemize}
\end{itemize}

\medskip

\subsection{National Payment Systems \& Switches of Bangladesh}

Bangladesh Bank operates and regulates several centralized payment platforms:

\begin{enumerate}
  \item \textbf{BACH (Bangladesh Automated Clearing House):}
  \begin{itemize}
    \item \textbf{BACPS (Bangladesh Automated Cheque Processing System):} Digital cheque truncation system where physical cheques are scanned and cleared through digital image exchange (High Value: $\ge$ Tk. 5 Lakh cleared within same day; Regular Value: cleared overnight).
    \item \textbf{BEFTN (Bangladesh Electronic Funds Transfer Network):} Batch-processing network for interbank electronic debit and credit transfers (salaries, pensions, dividends, vendor payments). Non-instant, processed in clearing cycles.
  \end{itemize}
  \item \textbf{RTGS (Real-Time Gross Settlement):}
  \begin{itemize}
    \item Designed for \textbf{high-value and urgent transactions} ($\ge$ Tk. 1 Lakh for individuals, unlimited for corporate).
    \item Settles on a \textbf{gross basis} (transaction by transaction without netting) in \textbf{real time} across central bank settlement accounts.
    \item Supports local currency (BDT) as well as selected foreign currencies (USD, EUR, GBP, CAD, JPY).
  \end{itemize}
  \item \textbf{NPSB (National Payment Switch Bangladesh):}
  \begin{itemize}
    \item Interoperable national routing switch connecting all bank ATMs, Point-of-Sale (POS) terminals, and \textbf{IBFT (Internet Banking Fund Transfer)}.
    \item Enables a cardholder of one bank to withdraw cash from any ATM or transfer funds instantly to another bank's account or card 24/7.
  \end{itemize}
  \item \textbf{Binimoy (IDTP --- Interoperable Digital Transaction Platform):}
  \begin{itemize}
    \item Launched by the ICT Division and Bangladesh Bank to enable interoperability between Banks, MFS providers (bKash, Rocket, Nagad), and Payment Service Providers (PSPs).
    \item Uses a Virtual Payment Address (VID) instead of exposing raw bank account or mobile numbers.
  \end{itemize}
  \item \textbf{SWIFT (Society for Worldwide Interbank Financial Telecommunication):}
  \begin{itemize}
    \item Secure messaging network connecting thousands of financial institutions globally for international wires, letters of credit, and treasury confirmations.
    \item Uses an 8- or 11-character \textbf{BIC (Bank Identifier Code)} or SWIFT Code (e.g., Sonali Bank Head Office BIC: \texttt{BSONBDDH}).
    \item Messages transition from legacy MT formats (e.g., \textbf{MT103} for single customer credit transfer, \textbf{MT202} for interbank transfer) to modern \textbf{ISO 20022 XML (MX)} standards.
  \end{itemize}
\end{enumerate}

\medskip

\subsection{Cybersecurity in Banking \& Lessons from 2016 BB Heist}

\begin{itemize}
  \item \textbf{Bangladesh Bank Cyber Heist (February 2016):}
  \begin{itemize}
    \item \textbf{Root Cause:} Attackers breached the internal network using phishing emails, infected workstations, and escalated privileges to access the \textbf{SWIFT Alliance Access (SAA)} terminal.
    \item \textbf{Vulnerabilities Exploited:} Lack of physical and logical network segregation (unmanaged switches, absence of dedicated VLANs separating SWIFT servers from general LAN), hardcoded credentials, and disabled printer monitoring.
    \item \textbf{Execution:} Attackers submitted 35 fraudulent SWIFT transfer requests totaling \$951 million to the Federal Reserve Bank of New York. \$81 million was routed to RCBC Bank in the Philippines before an anomalous typographical error (``fandation'' instead of ``foundation'') halted subsequent transfers.
  \end{itemize}
  \item \textbf{Post-2016 Regulatory Reforms \& Modern Banking Controls:}
  \begin{enumerate}
    \item \textbf{Air-Gapped / Isolated SWIFT Environment:} SWIFT infrastructure must reside in a dedicated demilitarized zone (DMZ) with multi-factor authentication, isolated from the general corporate network.
    \item \textbf{Hardware Security Module (HSM):} Tamper-resistant cryptographic hardware to securely store root digital certificates, PIN encryption keys, and transaction signing keys.
    \item \textbf{Security Operations Center (SOC) \& SIEM:} 24/7 continuous log aggregation, behavior anomaly detection, and automated endpoint detection and response (EDR).
    \item \textbf{Dual Authorization (Four-Eyes Principle):} Every financial data manipulation or high-value fund transfer must require two independent operators (Maker-Checker model).
    \item \textbf{PCI-DSS Compliance:} Payment Card Industry Data Security Standard certification required for processing, storing, or transmitting cardholder PINs and PAN numbers.
  \end{enumerate}
\end{itemize}

\medskip

\subsection{Anti-Money Laundering (AML) \& Regulatory Technology}

\begin{itemize}
  \item \textbf{BFIU (Bangladesh Financial Intelligence Unit):} The national central agency responsible for analyzing Suspicious Transaction Reports (STR) and Cash Transaction Reports (CTR) under the Money Laundering Prevention Act, 2012, and Anti-Terrorism Act, 2009.
  \item \textbf{Cash Transaction Report (CTR):} Commercial banks are mandated to submit a CTR to BFIU for any single cash deposit or withdrawal transaction of \textbf{Tk. 10 Lakh or above} (or foreign currency equivalent) within a calendar month.
  \item \textbf{Suspicious Transaction Report (STR) / SAR:} Any transaction that deviates from the customer's declared profile, lacks economic justification, or appears linked to illegal acts must be reported to BFIU promptly, regardless of amount.
  \item \textbf{KYC (Know Your Customer) \& e-KYC:} Mandatory identification procedure verifying customer identity, source of funds, and residential address. \textbf{e-KYC} integrates biometric fingerprint/facial scans with the national election commission database via the \textbf{Porichoy API} for instant paperless account opening.
  \item \textbf{Sanction Screening Engines:} Real-time automated cross-checking of customer names against United Nations Security Council (UNSC) resolution 1267 sanctions lists, OFAC, and local law enforcement terrorist databases prior to transaction execution.
\end{itemize}

\medskip

\section{High-Yield Banking IT Viva Questions \& Answers}

\subsection{Why should an IT graduate join a bank instead of a software development firm?}

\begin{quote}
\textit{Model Answer:}
``In a typical software firm, engineers often build isolated modules or consumer apps where downtime is inconvenient but rarely catastrophic. In contrast, a state-owned bank like Sonali Bank operates a mission-critical financial backbone. Here, software engineering and infrastructure directly safeguard national economic security, process billions of takas daily, and require six-nines (99.9999\%) availability. 

Working in banking IT exposes an engineer to enterprise-scale distributed systems, high-concurrency database locking, hardware cryptography, and stringent regulatory compliance. Additionally, it offers lifelong job stability, institutional prestige, and the opportunity to lead public-sector digital transformation.''
\end{quote}

\medskip

\subsection{What happens behind the scenes when a cardholder withdraws money from an ATM?}

\begin{enumerate}
  \item \textbf{Card Read \& PIN Encryption:} The ATM card reader captures the EMV chip data; the user enters their PIN on an Encrypted PIN Pad (EPP). The PIN is immediately encrypted into a secure PIN block using 3DES/AES keys managed by an onboard Hardware Security Module (HSM).
  \item \textbf{Message Formulation:} The ATM controller packages an \textbf{ISO 8583} financial transaction request message (containing card number, transaction type, withdrawal amount, and encrypted PIN block) and transmits it over an encrypted VPN to the bank's ATM Switch.
  \item \textbf{Routing via NPSB:} If the card belongs to another bank, the ATM Switch routes the ISO 8583 message to the \textbf{National Payment Switch Bangladesh (NPSB)}, which forwards it to the cardholder's Issuing Bank.
  \item \textbf{Authorization:} The Issuing Bank checks PIN validity with its HSM, verifies account status, checks balance and daily withdrawal limits, locks the requested funds, and sends an authorization approval code back along the return path.
  \item \textbf{Cash Dispense \& Settlement:} The local ATM dispenses the cash, generates a customer receipt, and sends a completion message back to the switch to debit the ledger permanently. If cash jams or the dispenser fails, a reverse transaction message cancels the provisional debit.
\end{enumerate}

\medskip

\subsection{How do you prevent insider threats and unauthorized data alterations in a bank's database?}

\begin{enumerate}
  \item \textbf{Privileged Access Management (PAM):} Database administrators (DBAs) should never have direct, unmonitored production access. Access must occur through bastions with session recording and temporary just-in-time credential issuance.
  \item \textbf{Separation of Duties \& Maker-Checker Model:} Application developers must not have write permissions to production databases. Any manual database correction must follow an automated change-management ticket signed off by risk officers.
  \item \textbf{Database Activity Monitoring (DAM):} Continuous auditing of all executed SQL statements, with automated alerts triggered if bulk update/delete statements or unindexed schema alterations occur.
  \item \textbf{Immutable Audit Logs \& Cryptographic Hashing:} Financial transaction logs should be replicated to write-once-read-many (WORM) storage or append-only log lakes with cryptographic hash chains to prevent retroactive tampering.
\end{enumerate}

\medskip

\subsection{What is the difference between BEFTN, RTGS, and NPSB in terms of speed, amount, and use cases?}

\begin{table}[htbp]
\centering
\small
\renewcommand{\arraystretch}{1.3}
\begin{tabular}{|l|p{3.8cm}|p{3.8cm}|p{4.2cm}|}
\hline
\textbf{Feature} & \textbf{BEFTN} & \textbf{RTGS} & \textbf{NPSB (IBFT / ATM / POS)} \\
\hline
\textbf{Settlement Type} & Net settlement in clearing cycles (Batch). & Gross settlement in real time (Transaction by transaction). & Real-time individual transaction switching. \\
\hline
\textbf{Speed} & Non-instant (Takes hours to next business day). & Instantaneous (within seconds to minutes). & Instantaneous (real-time 24/7/365). \\
\hline
\textbf{Value Limit} & Any amount (usually corporate payroll and dividends). & High value ($\ge$ Tk. 1,00,000 for individual; no upper limit). & Retail amounts (typically capped up to Tk. 3,00,000 to 5,00,000 per day). \\
\hline
\textbf{Currency} & Local currency (BDT) only. & BDT and approved foreign currencies (USD, EUR, GBP, etc.). & BDT only. \\
\hline
\textbf{Primary Use} & Bulk payroll, vendor checks, pensions, dividends. & Large interbank, corporate, and real-estate settlements. & ATM withdrawals, POS card swipes, instant app-to-app fund transfers. \\
\hline
\end{tabular}
\end{table}

\medskip

\subsection{What is your perspective on Central Bank Digital Currency (CBDC) for Bangladesh?}

\begin{quote}
\textit{Model Answer:}
``A Central Bank Digital Currency (a `Digital Taka') would represent a direct sovereign digital liability of Bangladesh Bank, unlike commercial bank money or private MFS wallets. 

While Bangladesh has made remarkable progress through MFS (bKash, Nagad) and agent banking, a retail CBDC could lower physical currency printing and logistics costs, facilitate programmable welfare distribution, reduce settlement latency in cross-border trade, and enhance financial transparency. However, its rollout must be carefully calibrated to avoid disintermediating commercial bank deposits, ensure resilient offline payment capability during network outages, and enforce military-grade cybersecurity protocols against sovereign cyber threats.''
\end{quote}
